\chapter{Wstęp}
\label{ch:wstep}

Dynamiczny rozwój technologii oraz rosnące wymagania stawiane współczesnym systemom technicznym powodują konieczność ciągłego doskonalenia istniejących rozwiązań oraz poszukiwania nowych metod ich optymalizacji. W szczególności dotyczy to obszaru [tu wstaw dziedzinę, np. automatyki, informatyki, elektroniki, mechaniki], w którym innowacje mają bezpośredni wpływ na efektywność, bezpieczeństwo oraz komfort użytkowania systemów.
Celem niniejszej pracy inżynierskiej jest [tu wpisz cel pracy, np. zaprojektowanie, analiza, implementacja, optymalizacja] rozwiązania z zakresu [obszar tematyczny]. Realizacja tego celu wymagała zarówno przeglądu aktualnego stanu wiedzy, jak i zastosowania odpowiednich narzędzi analitycznych oraz projektowych.
Zakres pracy obejmuje [krótki opis zakresu, np. analizę literatury, opracowanie modelu, wykonanie projektu, implementację systemu, przeprowadzenie testów]. W trakcie realizacji projektu wykorzystano [narzędzia / technologie / metody badawcze], które umożliwiły osiągnięcie założonych rezultatów.
Praca została podzielona na kilka rozdziałów. W pierwszym z nich przedstawiono podstawy teoretyczne oraz przegląd istniejących rozwiązań. Kolejne rozdziały zawierają opis opracowanej koncepcji, proces projektowania oraz wyniki przeprowadzonych badań i testów. Całość kończy podsumowanie wraz z wnioskami oraz propozycjami dalszych kierunków rozwoju.